\section{Discussion}
\label{sec:discussion}

% Status: partial DRAFT.  Only the evidence-supported themes are written.
% Everything depending on Phase 8, the calibration ladder, or the pool
% campaign is marked and must not be written before the data exist.

\subsection{Model Complexity Does Not Imply Better Prediction}

The estimator ablation supports a claim that is narrower and more useful than
the one usually made for Kalman filtering in tracking pipelines. Adding a
constant-velocity motion model did not improve prediction where prediction
was hardest. During the dropout at maneuver onset the hold was roughly twice
as accurate as the constant-velocity filter, and the reason is visible in the
reference signal rather than inferred: the vehicle's own commitment
introduces a step in the obstacle's image-space velocity, so the filter
extrapolates a velocity that has just ceased to hold, while the hold's
implicit zero-velocity assumption happens to be closer to the truth
immediately after the step. The Kalman filter nonetheless produced the only
clean outlier scenario, because $\chi^2$ gating rejects a corrupted
measurement that the smoothing variant partially absorbs.

The generalizable point is that prediction accuracy and outlier robustness
are separate benefits that are routinely conflated. A filter can be the
better choice for one and the worse choice for the other, and reporting a
single aggregate success rate hides the trade. For the present study this has
a concrete methodological consequence: the two benefits are stressed by
different physical phenomena --- ego-motion dynamics and detector failure
modes respectively --- so a simulator could reproduce one faithfully while
misrepresenting the other, and the calibration ladder is instrumented to tell
those cases apart.

\subsection{What an Ego-Motion-Aware Model Would Have to Do}

The mechanism above suggests, but does not establish, that the useful
correction is to inform the motion model about the vehicle's own commanded
maneuver. An exploratory damped variant sat between the hold and the
constant-velocity filter without beating the hold, which is consistent with
damping addressing the symptom rather than the cause. We flag this as a
hypothesis: a model that predicts image-space motion from the commanded body
velocity should dominate both, and testing it is future work rather than a
result of this paper.

\subsection{Fidelity Is Not Predictivity}

The framing of this paper rests on a distinction that the underwater
literature currently under-uses. High simulator fidelity is a property of the
model; predictivity is a property of the model \emph{as an evaluation
instrument}, and the second does not follow from the first. Truong et
al.~\cite{truong2022rethinki} demonstrated the gap in the training regime,
and Kadian et al.~\cite{kadian2020simrealp} demonstrated that predictivity
can be improved deliberately in the evaluation regime. Underwater work has
concentrated on the first property --- richer hydrodynamics, richer sensor
models, photogrammetric
scenes~\cite{benzon2022anopenso,orjales2025towardsa,mari2026deeprein} --- and
reports transfer success as evidence of it. The claim that our ladder
improves predictivity is not yet supported by data and is deliberately not
made here; what this paper fixes is the instrumentation that can decide it.
\SREAL

\subsection{Planner Comparison}

% TODO_PHASE8_FINAL_RESULT: written only after the Phase-8 campaign and the
% D-013 freeze.  Do not anticipate the outcome, in either direction.
\PEIGHT
